\chapter*{Start}\label{start}
\addcontentsline{toc}{chapter}{Start}

\enquote{\noindent\textsc{{How many kids would you have had?}} I mean, if it were not for\dots{}}

Her voice slowly fades out, her face shying away the same way she does every time the topic of her mother is approached. So many times, so many years, and it still pains me to see how she gets the guilty blow-back from herself every time she does it, and yet she cannot avoid repeating it.

I say, \enquote{Hey Ladybug, don't get blue---it does not hurt like before.} It's a lie, it still hurts. And she senses it somehow.

I reach out for her, bringing her to a hug as she shyly and yet dramatically trudges into the living room from the apartment's entrance. I sniff her hair for a while and notice: smoke. I ask, \enquote{Were you smoking?}

She takes some distance, opens her mouth demonstratively and exhales in my direction: negative for smoke. 

\enquote{It was Justyna,} she says. \enquote{We gathered in the park, she's got the confirmation of a position. She will start right after finishing training, in some months. She was musing that, with her and Thomas having confirmed positions, they can start to make plans for a baby.}

\enquote{So, \emph{this} is where this kid-number-story comes from,} I comment. \enquote{Now tell me: how will Justyna have a baby and smoke? Or\dots{}}

\enquote{\dots~she will have to stop when she decides to get pregnant, so she wants to smoke now while she can,} she complements my line of thought. She is good at doing that---just like her mother. I usually hate when people try to predict what others are about to say, but only because most of the time people get it wrong. 

Nitzanah never does it. 

I invite her, with gestures, to sit down. She does it, drops her backpack near the couch. She still dresses like a kid, even though she's a woman old enough to muse about having kids in earnest. I feel glad that I enjoyed her years as a kid as fully as I could. 

I decide that it starts now. 

\enquote{Nitzy\dots{}} She understands by the undertone that her attention is required.

\enquote{Yes?}

\enquote{Did Justyna comment on pre-natal precautions, like genetic screening?} I ask.

\enquote{No\dots~is that not only recommended for persons at risk of some problem?}

I answer, \enquote{Mostly, yes---and I do not know if she or her partner---Thomas?---are at risk\dots~but perhaps it is important that you know about us.}

Her whole posture and attitude changes, now she's straightened up and looking at me very earnestly. She's waiting for the continuation. I abide.

\enquote{For one, I do not believe I told you before, but I am a thalassemia carrier. It is a genetic condition that can cause anemia of diverse degrees, depending on how many genes are affected. Now, you do not have anemia, and probably do not have any gene affected, but there are several conditions that may run silent and still be a problem for the children. And there are things about your mother's side you do not know, but should, before you have children.}

She looks at me for a while. After some reflection that I recognize from oh so many sweet moments, she declares, \enquote{You are preparing me to get away from you. To be alone. Stop.}

I laugh out loud and reply \enquote{Of course I am not! Do not be so dramatic. Just like your mother, aren't you? Just look at you: a fully grown woman. If Justyna can start having kids anytime, so do you. And if you do, why not check things beforehand?}

She smiles at me in a warm way. It breaks my heart not being capable of telling the truth as it is and hiding behind a lab result.